\documentclass[11pt,a4paper]{article}

\usepackage{amsmath,amssymb,amsfonts}
\usepackage{geometry}
\usepackage{hyperref}
\usepackage{booktabs}
\usepackage{xcolor}

\geometry{margin=1in}
\hypersetup{colorlinks=true, linkcolor=blue, citecolor=blue, urlcolor=blue}

\newcommand{\workingnote}[1]{\textcolor{red}{\textbf{[v0.1 NOTE: #1]}}}

\title{\textbf{A Finite Internal Dimension from the Condensate:\\
Deconstruction, the Composite Spectrum,\\
and a Falsifiability Test}\\[1ex]
\large Working note v0.22 --- NOT a manuscript.\\
Status: concept + test design + partial numerics; v0.20 records a
MEASUREMENT-FORCED CORRECTION: the induced Skyrme coefficient is
now known exactly and is NEGATIVE
($\kappa_U = -17/(1152\pi^2)$, chirally complete, closed form +
Pauli--Villars confirmed), retracting this note's earlier
positive-sign entries (S2-1, Gate 3) and re-scoping soliton
stabilization to the endogenous flavor-singlet $\omega$ channel
(existence and dynamization gates passed; continuum-radius gate
open).
v0.22 records the consequences of that same measurement for gate
S2-2: a $3{+}1$D physics correction inherited from the companion
dark-energy paper, the monopole-action arithmetic
$S_{\rm mono} = -85/(9\pi^2)$, and the retraction of the
Coulomb-side resolution that the superseded $\kappa$ numbers had
supported; and it states the provenance status of the $\omega$
gates.}

\author{Zeta Hoi-Ho Cheng \\
\textit{Independent Researcher, Ontario, Canada}}

\date{July 2026 (v0.22)}

\begin{document}
\maketitle

\begin{abstract}
We record a hypothesis and its kill criteria.
The hypothesis: the composite spectrum of the $H(4)$ lattice
fermion condensate is organized by an emergent finite internal
dimension, in the precise sense of dimensional deconstruction ---
$N$ four-dimensional fields with a specific hopping mass matrix
are the same theory as one field on $M_4$ times an $N$-site
internal lattice, and the condensate's Goldstone modes are
candidates for the link fields that build the internal hops.
On this reading, the composite masses of the programme
($m_P$, $m_\sigma$, $m_V$, $m_A$) are not independent parameters
but must fall, as \emph{ratios}, on the eigenvalue spectrum of a
discrete one-dimensional Laplacian for some small $N$ and boundary
condition.
This note fixes the dictionary, designs the spectral test, records
the numerical results in hand (updated in v0.19 to the overlap
frame: an exact Goldstone $m_\pi = 0$, a measured $m_\sigma$ law
with chiral-limit intercept unity, a unified $S/P$ criticality and
a degenerate $V/A$ criticality map---with the caveat that at the
invariant-universal coupling the vector channel supports no
below-threshold bound state; the earlier Wilson-frame table, with
its blocked scalar channel and axial floor, is retained as the
historical record of the diagnosis), and
states what would kill the hypothesis.
Two of this note's own entries have since been retracted by
measurement and are recorded as such: the induced quartic-gradient
(Skyrme) coefficient, previously reported positive, is now known
exactly and is \emph{negative},
$\kappa_U = -17/(1152\pi^2)$, which withdraws both the one-loop
soliton stabilizer and the positive Berry--Maxwell term inherited
by the orientation layer's photon candidate (gate S2-1), and
suspends the monopole gate S2-2 behind it; soliton stabilization
is re-scoped onto the endogenous flavor-singlet $\omega$ channel,
whose continuum-radius gate is open.
No claim of confirmation is made or available at this stage.
\end{abstract}

% ===================================================================
\section{Scope and non-claims}
\label{sec:scope}
% ===================================================================

This note is deliberately narrower than the intuition that
motivated it.
We do \emph{not} propose an interpretation of quantum mechanics;
no claim is made that quantum uncertainty, the Born rule, or
exchange statistics originate in motion along a hidden dimension,
and nothing here engages Bell-type theorems, because nothing here
needs to: deconstruction reorganizes \emph{field content}, taking
quantum field theory as given.
The only physical question asked is structural:

\begin{quote}
\emph{Do the composite masses of the condensate programme fit the
eigenvalue spectrum of a discrete internal Laplacian for some
small $N$, so that their ratios become outputs rather than
inputs?}
\end{quote}

If yes, the internal-dimension language earns its keep as an
organizing principle and parameter-reducer.
If no, the language is decoration for this model and the
hypothesis is dead in its minimal form.
Both outcomes are recorded as acceptable.

% ===================================================================
\section{The dictionary: deconstruction in brief}
\label{sec:dict}
% ===================================================================

\subsection{One field on an internal lattice $\equiv$ $N$ fields
with a hopping matrix}

Let $\phi(x,n)$, $n = 0,\dots,N-1$, live on $M_4$ times an
$N$-site internal lattice with spacing $b$ and hopping term
$\frac{1}{b^2}\sum_n |\phi(x,n{+}1)-\phi(x,n)|^2$.
Discrete Fourier transform in $n$ diagonalizes the hopping matrix;
for a periodic ring the four-dimensional mass tower is
\begin{equation}
m_k^2 = \frac{4}{b^2}\,\sin^2\!\frac{\pi k}{N},
\qquad k = 0,1,\dots,N-1,
\label{eq:ring}
\end{equation}
while an open chain (free boundaries) gives the cosine-type
spectrum
\begin{equation}
m_k^2 = \frac{4}{b^2}\,\sin^2\!\frac{\pi k}{2N},
\qquad k = 0,1,\dots,N-1,
\label{eq:chain}
\end{equation}
and fixed or twisted boundaries shift the allowed $k$.
Small $k$ reproduces a Kaluza--Klein tower with
$1/R \to 1/(Nb)$.
The converse statement is the content of dimensional
deconstruction (Arkani-Hamed--Cohen--Georgi; Hill--Pokorski--Wang,
2001): $N$ four-dimensional fields whose mass matrix is a discrete
Laplacian \emph{are} a lattice extra dimension.
Whether ``the dimension exists'' is therefore not an ontological
question; the only physical content is whether the spectrum
follows a discrete-Laplacian eigenvalue structure.

\subsection{Where the links come from: condensate Goldstones}

In the original moose construction the link fields connecting
adjacent sites are not fundamental: they are nonlinear
sigma-model fields --- Goldstone modes of fermion condensates.
The arrow of explanation is
\begin{equation}
\text{condensate} \;\to\; \text{Goldstone links} \;\to\;
\text{hopping matrix} \;\to\; \text{internal dimension},
\label{eq:arrow}
\end{equation}
which matches the direction of the present programme: the
dimension, if present, is an \emph{output} of the condensate's
internal structure, not a postulated background.
The angular pseudo-Goldstone of the complex condensate
(Papers~2 and~4) is the first candidate link degree of freedom;
whether the $H(4)$ condensate's internal structure is rich enough
to furnish a link \emph{chain} (rather than a single phase) is an
open structural question and part of the Fierz analysis
(Section~\ref{sec:gates}).

\subsection{Fermion side: domain-wall correspondence}

The fermionic counterpart of a finite internal dimension with a
mass profile is the domain-wall construction: chiral zero modes
localize on the boundaries, and four-dimensional chirality is
reread as boundary localization in the internal direction.
This is the rigorous residue of the original ``spin from a fourth
dimension'' intuition --- weaker, but exact.
It is deliberately \emph{not} the entry point of the test
(no independent fermionic spectrum data exist); it is the second
step, to be revisited only if the bosonic spectral test passes.

% ===================================================================
\section{The spectral test}
\label{sec:test}
% ===================================================================

\subsection{Design}

The test operates on composite-boson mass \emph{ratios} at fixed
microscopic couplings:
\begin{enumerate}
\item Assemble the composite spectrum
$\{m_P,\, m_\sigma,\, m_V,\, m_A\}$ in lattice units at a single
$(G, m_f)$ point, with couplings across channels fixed by the
Fierz decomposition of the microscopic four-fermion structure
(not tuned independently).
\item Ask whether the mass-squared ratios coincide, within
combined errors, with eigenvalue ratios of
Eq.~\eqref{eq:ring} or Eq.~\eqref{eq:chain} (or a twisted variant)
for some $N \le 5$.
The overall scale $4/b^2$ is one free parameter; the ratios are
not.
\item The pseudoscalar, being Goldstone-adjacent, is a natural
$k=0$ zero mode; a spectrum with an exact (or near) zero mode
favors periodic or free boundaries.
\end{enumerate}

\subsection{Kill criteria}

The hypothesis is dead in its minimal form if any of the
following holds:
(i) no $N \le 5$ and boundary condition reproduces the measured
ratios within errors;
(ii) the Fierz-fixed couplings, once computed, remove the freedom
needed to place even two states consistently on any candidate
spectrum;
(iii) the chirally clean scalar extraction places $m_\sigma$ in a
position incompatible with every candidate spectrum that
accommodates $m_V$ and the $m_A$ floor.
A fit achieved only by letting $N$ grow or by tuning couplings
per channel is a failure, not a success: the entire value of the
hypothesis is parameter reduction.

% ===================================================================
\section{Partial results in hand (Wilson-frame run, July 2026;
superseded in part --- see the update at the end of this section)}
\label{sec:partial}
% ===================================================================

All numbers below are from the calibrated bubble pipeline
(Wilson $r=1$, volumes to $160^4$, conventions anchored to the
companion vector-channel paper: published $\Pi_V(0)$ at
$m_f a = 0.2, 0.5$ reproduced to better than $0.1\%$;
$\Pi_A(0) \simeq -0.19$ reproduced; the gap-equation input
$I_0 = 0.0844$ of the gravity paper independently reproduced).
EFT composite masses are defined by
$m_\Gamma^2 = [\,1/G_\Gamma - \Pi_\Gamma(0)\,]/Z_\Gamma$ with
$Z_\Gamma = d\Pi_\Gamma/dE^2$ from the analytic continuation
$q_0 \to iE$ of the finite momentum sum.

\begin{center}
\begin{tabular}{lccl}
\toprule
Channel & $\Pi_\Gamma(0)$ ($m_f a{=}0.05$) & $Z_\Gamma$ & Status \\
\midrule
$P$ & $+0.331$ & $+0.081$ & Goldstone-adjacent; candidate $k{=}0$ \\
$V$ (spatial) & $+0.284$ & $+0.058$ & healthy; mass dial
  $m_V^2 a^2 = [1/G_V - 0.284]/0.058$ \\
$A$ (spatial) & $-0.189$ & $+0.044$ & \emph{floor}:
  $m_A a \ge 2.08$ even at $G_A \to \infty$ \\
$S$ & $+0.141$ & $-0.066$ & blocked: Wilson-contaminated vertex,
  wrong-sign $Z$ \\
\bottomrule
\end{tabular}
\end{center}

Notes.
(1) The axial floor is the first quantitative form of ``axial
permanently heavy'' (the companion paper had the sign argument
only); it is an extrapolation of the $q^2$-expansion beyond its
small-$E$ regime and is indicative.
At $m_f a = 0.2$ the floor is $m_A a \ge 3.18$.
(2) A consistency check: the vector-paper scenario
$m_V = e^{-2}\Lambda$ with $a\Lambda = \pi$ requires
$G_V = 3.40$, at the lower edge of the published criticality band
$G_c^V = 3.4$--$4.4$.
(3) The scalar block is understood, not merely observed: the
lattice $\bar\psi\psi$ bubble fails the continuum consistency
condition (shared-coupling Goldstone $+$ $\sigma$ pole at $2m_f$),
its induced kinetic term carries the wrong sign at small momenta,
and a central-Brillouin-zone mask does not repair it because the
Wilson mass is $\mathcal{O}(1)$ across the zone.
The same pathology class blocked the direct $\beta_F$ extraction
in the gravity paper.
(4) The gap-versus-bubble criticality spread
($3.02 / 5.93 / 7.09$) is documented in the gravity paper v2.9
and caps what scalar-channel bubble numbers may be used for.

\emph{Test status at the time of the Wilson run: not yet runnable.}
Hard data points: $m_P \approx 0$ and the $m_A$ floor.
Dial: $m_V(G_V)$, awaiting Fierz.
Missing: $m_\sigma$ (blocked), Fierz-fixed coupling ratios.

\subsection*{Overlap-frame update (v0.19): the table above is
superseded}

The Ginsparg--Wilson (overlap) campaign has since replaced the
Wilson-frame numbers above; the table is retained as the
historical record of the diagnosis (the Wilson scalar block, in
particular, is confirmed as a discretization artifact).  The
current spectral inputs are:
\begin{itemize}
\item \textbf{$S/P$ sector unified and healthy.}  Exact
$\Pi_S = \Pi_P$ degeneracy by the lattice chiral Ward identity;
single criticality $G_c^{S/P} = 1/0.130 = 7.69$ (overlap units,
$\rho = 1$); $Z_S = +0.065 > 0$---the scalar channel is no longer
blocked.
\item \textbf{$m_\sigma$ measured.}  True-pole scan across
$m_{\rm dyn} a = 0.1$--$0.4$:
$m_\sigma/(2m_{\rm dyn}) = 1 - 0.98\,(m_{\rm dyn}a)
+ 0.72\,(m_{\rm dyn}a)^2$, chiral-limit intercept $1.000(10)$
(details in the spectral-test anchor of
Section~\ref{sec:u3seed}).
\item \textbf{$V/A$ sector.}  With chirally improved vertices,
exactly degenerate at $m_f \to 0$: $G_c^{V/A} = 1/0.065 = 15.4$;
$Z_V > 0$.  The Wilson-frame axial floor is retired in this frame
(the $V$--$A$ splitting grows with $m_f$ as chiral symmetry
predicts, rather than the axial being permanently heavy).
\item \textbf{Tensor.}  $T_T \to 0$ in the chiral limit
(continuum-consistent); $G_c^T \to \infty$.
\item \textbf{Vector-state caveat.}  The companion vector-channel
paper's pole-level analysis shows that at the invariant-universal
coupling $G = G_\chi = 0.5\,G_c^{V/A}$ the vector channel supports
\emph{no below-threshold bound state}; a genuine light vector
requires sub-percent proximity to $G_c^{V/A}$.  Any candidate
tower must therefore either place the vector on the continuum
(resonance) side or restrict itself to the states that exist.
\end{itemize}
\emph{Test status now: the mass side is anchored}
($m_\pi = 0$ exact, the $m_\sigma$ law above, and the $V/A$
criticality map with the vector-state caveat); \emph{the coupling
gate remains the outstanding requirement}, in re-scoped form: the
minimal chiral pair generates exactly zero $V/A$ coupling by Fierz
(machine-verified), so cross-channel coupling ratios are not
Fierz-inherited; in the adopted invariant-universal model they are
fixed instead by the single-coupling economy postulate
($G_V = G_A = G_\chi$).  Whether that postulate, combined with the
vector-state caveat, leaves enough of a tower for the $N \le 5$
fit to be meaningful is an open design question recorded in
Section~\ref{sec:gates}.


% ===================================================================
\section{Near-critical geometry: the equivalence-principle test}
\label{sec:epgeom}
% ===================================================================

\subsection{The conjecture, and the death of its naive form}

The intuition: near-criticality is not parameter tuning but the
mathematical condition for a discrete substrate to admit a
continuum geometry---when $\xi \gg a$, excitations cannot resolve
the lattice and hypercubic anisotropies wash out.
This is Wilsonian renormalization-group lore, and its standard
quantitative form is \emph{power-law} irrelevance:
species-dependent metric responses suppressed as $(a/\xi)^p$,
$p > 0$.

\textbf{The power-law form is refuted by our own measurement.}
A species-resolved light-cone test (uniform temporal deformation
of the substrate; effective speeds $c_\Gamma^2 = Z^{(s)}_\Gamma /
Z^{(t)}_\Gamma$ per channel; built-in isotropy validation at
$\varepsilon = 0$) gives, across three fermion masses
$m_f a = 0.2, 0.1, 0.05$, a composite--fermion cone split of
$10\% \to 7.1\% \to 5.6\%$ (pseudoscalar; vector similar), lying
on the one-parameter curve
\begin{equation}
\delta c^2_\Gamma \;=\; \frac{A_\Gamma}{\ln\!\bigl(1/(m_f a)\bigr)},
\qquad A_P \simeq 0.16,\; A_V \simeq 0.19,
\label{eq:loglaw}
\end{equation}
to within $0.3$ percentage points at every point:
\emph{logarithmic}, not power-law, restoration.
The mechanism is identified: $c^2$ is a \emph{matching
coefficient} fixed where the composite is assembled, and the
bound-state wavefunction has $\mathcal{O}(1)$ support at the
cutoff, so ultraviolet anisotropy leaks in at $1/\log$ rather
than being power-suppressed.  Wilsonian irrelevance protects the
effective operators of an already-light field; it does not
protect compositeness matching.  This is precisely why the
Collins--Perez--Sudarsky objection to emergent gravity is
dangerous, now exhibited quantitatively within the model.

Verdict: near-critical protection \emph{exists} (leading-log
intercept zero; the near-critical limit is the continuum-geometry
limit), but in the weakened logarithmic form of
Eq.~\eqref{eq:loglaw}.

\subsection{The phenomenological budget, honestly}

Extrapolating Eq.~\eqref{eq:loglaw} to physical scales,
$\delta c^2 \sim A/\ln(M_{\rm Pl}/m) \sim 10^{-3}$ even for
ultralight composites, against multi-messenger and
E\"otv\"os-class bounds at $10^{-15}$: twelve orders short.
Logarithmic protection alone is phenomenologically insufficient.
Survival requires a second layer, and the test data contain its
candidate: the \emph{composite--composite} split.
\textbf{Correction (v0.4):} the v0.3 claim that this split falls
faster than $1/L$ is \emph{retracted}: it rested on
volume-unconverged points.  A dedicated audit using an
exact-difference reformulation ($N_P - N_{V_2} = 8 s_2 s_2'$,
removing large-number cancellation) revealed strong, slowly
saturating finite-volume dependence of the split
(at $m_f a = 0.1$: $\sigma = 0.0064, 0.0173, 0.0253, 0.0265$
for $n = 48, 64, 96, 128$).  The saturated values are
$\sigma_\infty(0.2) = 0.036(1)$ and
$\sigma_\infty(0.1) = 0.027(1)$, giving a per-halving ratio
$0.75(4)$---consistent with $1/L$ (predicts $0.70$) and
$\sqrt{m}$ ($0.71$), excluding $1/L^2$ and linear $m$ ($0.50$).
The composite--composite split therefore falls only
\emph{logarithmically}, like the composite--fermion split, with
$A_{PV} \simeq 0.06$; extrapolation to physical scales leaves
$\sim 4\times10^{-4}$, eleven orders above the $10^{-15}$
bounds.  \textbf{At the one-loop (free-bubble) level of the
present computation, leg~(2) of the trichotomy fails: no
second-layer power-law protection is visible.}
The remaining out is genuine but uncomputed: the present cones
are built from free-fermion bubbles, whereas near criticality the
composite sector is strongly interacting, and emergent-symmetry
mechanisms at interacting fixed points---precisely what one-loop
misses---are the known route by which irrelevant anisotropies can
acquire faster decay.  Whether the interacting theory rescues
leg~(2) is now the sharpest open question of the programme.

\subsection{Survival trichotomy (all three required)}

\begin{enumerate}
\item \textbf{Standard-Model matter must embed as composites.}
The observable cone comparisons (graviton vs.\ photon vs.\
matter) are then composite--composite; the substrate-fermion cone
drops out of observables.  The LIV analysis thereby upgrades
``SM as composites'' from an aesthetic choice to a
\emph{structural requirement} of the programme.
\item \textbf{Composite--composite splits must run below
$10^{-15}$} at observable scales.  This requires a dedicated
error-controlled measurement of the $P$--$V$ split scaling law:
same-footing speed definitions for all species, finite-difference
extrapolation, larger volumes, lighter masses.  This is the
highest-priority computation exposed by the test.
\item \textbf{The substrate fermion must be unobservable as a
propagating species} at accessible energies (dark substrate),
so that its $\mathcal{O}(10^{-3})$ cone mismatch is not an
observable LIV.
\end{enumerate}
Failure of any leg is falsification-grade for the programme.

\subsection{Relation to the internal-dimension hypothesis}

Two connections.  First, the conjecture of this section and the
internal-dimension hypothesis of Sections~\ref{sec:dict}%
--\ref{sec:test} are the same kind of statement: both assert
that structures invisible at the lattice scale (a continuum
geometry; an internal spectrum) organize the low-energy theory,
and both are given pre-registered quantitative forms that can
fail.  The naive power-law form of the present conjecture
\emph{did} fail; its surviving logarithmic form is the sharper
object.  Second, the intuition that lattice sites are not points
but carry extension (``the site is a line, not a point'') is,
in its controllable form, \emph{identical} to the
internal-dimension hypothesis: a site carrying a continuous
internal interval is the $N \to \infty$ limit of the $N$-layer
internal lattice of Section~\ref{sec:dict}, and its observable
signature is the same tower structure gated by the same spectral
test.  Claims about site extension that do not modify any
computable (spectrum, cone splits, couplings) are presentational
and carry no evidential weight.  One genuinely open question is
recorded, without any supporting computation: whether internal
site structure modifies the matching-coefficient mechanism and
hence the $1/\log$ law of Eq.~\eqref{eq:loglaw}; there is at
present no reason to expect improvement, since additional
ultraviolet structure generically adds to, rather than removes,
cutoff-scale anisotropy.


% ===================================================================
\section{Matter as topology: the $S^3$ extension and its gauntlet}
\label{sec:s3}
% ===================================================================

A topological audit of the minimal chiral pair closes one door and
opens a specific other.  The minimal model's vacuum manifold is the
chiral circle $S^1$, and $\pi_2(S^1) = \pi_3(S^1) = 0$:
\emph{point-like solitonic matter does not exist in the minimal
model}---its topological inventory contains only strings
($\pi_1 = \mathbb{Z}$) and, under the Wilson tilt, domain walls,
both cosmological liabilities rather than matter candidates.
If matter is to be a stable localized structure of the condensate,
the vacuum manifold must be enlarged; the minimal enlargement with
$\pi_3 = \mathbb{Z}$ is $S^3$, realized by the flavor-$SU(2)$
chiral pair
\begin{equation}
\mathcal{L}_{\rm int} = \frac{G}{2N}\Bigl[(\bar\psi\psi)^2
+ \sum_a (\bar\psi\, i\gamma_5\tau^a\psi)^2\Bigr],
\qquad (\sigma,\vec\pi) \in S^3,
\label{eq:su2pair}
\end{equation}
with skyrmionic matter (mass $=$ trapped energy, conserved charge
$=$ winding number, both induced from the fermion determinant as
in the Skyrme literature) and, as a bonus, $\pi_1(S^3) = 0$
removing the cosmic-string liability of the $S^1$ model.
This is the controlled form of the ``matter as a bubble of the
condensate'' intuition, and it links to the internal-structure
hypothesis of this note: topology \emph{requires} a condensate
with more internal structure than the minimal pair if the model is
to contain matter at all.

The extension was passed through the same machine-verified Fierz
gauntlet that excluded the universal coupling (particle--hole
exchange weights and exact pairing decomposition, residual
$10^{-14}$).  Verdicts:
(i) the topological gains stand and are $N$-independent;
(ii) the hope that flavor Fierz regenerates a vector coupling is
dead---the induced vector-singlet coupling is \emph{repulsive}
(exchange weight $+0.50$), so the $G_V$ orphan problem is not
solved here;
(iii) the $S^1$ model's pairing \emph{sign protection is lost}:
the axial-diquark (flavor singlet and $\tau_{1,3}$) and
tensor-diquark ($\tau_2$) coefficients turn positive
($+0.25$, $+0.50$) against positive pair bubbles.  The
normalizations are now pinned (machine-verified Wick factor $2$
for diagonal pair components; anomalous-equals-normal bubble
identity verified in the tensor channel to $10^{-4}$), giving the
criticality condition $N < N_{\rm crit} = 2\,G_{\rm op}\,c_i
T_i$ with the operating point at the $S^3$ condensate's own
criticality ($G_{\rm op} = G_c/2$ from the isospin doubling of
the tadpole).  Result: $N_{\rm crit} < 1$ in both the gap and
bubble conventions---the vacuum is \emph{magnitude-safe at every
$N \geq 1$}.  The thinnest margin is the tensor diquark at
$N = 1$: $1.4\times$ (gap convention) to $2.8\times$ (bubble
convention), where $1/N$ corrections are uncontrolled, so
$N \geq 2$ ($\geq 2.9\times$ everywhere) is the robust domain
and $N = 1$ is admissible but thin;
(iv) an attractive tensor-singlet particle--hole coupling appears
at $\mathcal{O}(G/N)$, subcritical by a factor $\gtrsim 3$ at any
$N \geq 1$;
(v) Eq.~\eqref{eq:su2pair} is not $U(1)_A$-symmetric, so the
dark-energy axion identification of the companion programme must
be re-audited in this extension (a $U(2)\times U(2)$ completion
would require its own gauntlet pass).
\textbf{Correction (v0.20), superseding the v0.18 entry.}  The
earlier ``positive, stable across masses and schemes''
$\kappa$ result is retracted: it was a fit-leakage artifact---the
box-only diagram set is chirally incomplete (nonvanishing
zero-momentum amplitude) and a fit basis lacking a constant term
leaked that amplitude into the $k^4$ coefficient, mimicking even
the expected chiral-log growth.  The corrected, chirally complete
measurement (five diagram classes; pointwise $A_0 \equiv 0$ at
machine precision; Ward-anchored normalization; five-point mass
scan; two evaluators, two integrators) is now closed in exact
form,
\begin{equation}
\kappa_U \;=\; -\frac{17}{1152\pi^2} \;=\; -0.0014952,
\end{equation}
confirmed under Pauli--Villars regularization and agreeing with
the Monte-Carlo value to $0.3\%$.  The sign is \emph{negative}:
the quartic-gradient term is anti-Derrick at one loop, and this
note's stabilization claims built on $\kappa > 0$ are withdrawn
below (S2-1, Gate 3).
\emph{The constructive replacement.}  Stabilization is
transferred to a channel this note's own Gate-0 audit already
recorded: the $U(3)$ interaction's particle--hole exchange
contains exactly one nonzero channel, a \emph{repulsive}
vector/axial singlet pair.  The Fierz identity is now pinned with
anchors and full-tensor reconstruction ($G_\omega = -G/N$ in the
companion convention: Walecka-repulsive, pure vector-singlet, the
flavor octet vanishing identically), and the channel has been
Hubbard--Stratonovich dynamized with all four health gates
passing: induced kinetic $Z_\omega > 0$; non-tachyonic
$M_\omega^2 = 1/(g_0 Z_\omega) > 0$; WZW/Goldstone--Wilczek--anchored topological-current coupling
$g_{\omega B} = Z_\omega^{-1/2} \neq 0$
($c_{\rm GW} = 1$ reproduced); and a static
kernel $D_{00}(0,q) > 0$ for all $q$---the repulsive branch's RPA
denominator $1 + g_0\Pi_V$ is positive definite, the exact mirror
of the attractive channel's critical denominator.
\emph{Provenance status (v0.22).}  Of these, only the
Goldstone--Wilczek normalization ($c_{\rm GW} = 1$) currently has
migrated, independently re-runnable artifacts in this paper's
repository; the $Z_\omega$, $M_\omega^2$, $D_{00}$ and Fierz
($G_\omega = -G/N$) results reside in a legacy branch whose scope
is mixed with the companion vector and dark-matter papers, and
await scope-separated migration and independent review.  They are
therefore recorded here as \emph{computed but not yet
independently certified}---a distinction this note takes
seriously, since the retraction above was itself caused by an
uncaught fit artifact in numbers that had looked stable.  In the \emph{local heavy-vector limit} the truncated landscape
$E(R) = AR - |B|/R + C/R^3$ possesses a unique stable minimum for
any $C > 0$; the full $\omega$ response is the nonlocal kernel
$E_\omega = \tfrac{g_{\omega B}^2}{2}\int B^0 D_{00} B^0$, and the
decisive, registered continuum gate is the full nonlocal radius
scan: $R_\star \gg a$ with $E''(R_\star) > 0$ (a cutoff-scale
minimum would be a lattice lump, not continuum topological
matter).  Soliton existence is \emph{not} claimed here; it is
re-scoped onto this $\omega$-stabilized landscape pending that
gate.
The final gate, the axion re-audit, is now closed---by
superseding Eq.~\eqref{eq:su2pair} rather than repairing it.
The dark-energy programme requires the $U(1)_A$ breaking to be
\emph{soft} (topological only), whereas Eq.~\eqref{eq:su2pair}
breaks it at tree level with $\mathcal{O}(G)$ strength, giving
the would-be axion a cutoff-scale mass.  The minimal cure is the
classically $U(2)\times U(2)$-symmetric interaction
\begin{equation}
\mathcal{L}_{\rm int} = \frac{G}{2N}\Bigl[(\bar\psi\psi)^2
+ (\bar\psi i\gamma_5\psi)^2
+ (\bar\psi\tau^a\psi)^2
+ (\bar\psi i\gamma_5\tau^a\psi)^2\Bigr],
\label{eq:u2pair}
\end{equation}
whose vacuum manifold is $U(2) \cong (S^1\times S^3)/\mathbb{Z}_2$:
the $S^3$ factor retains skyrmionic matter
($\pi_3(U(2)) = \mathbb{Z}$) while the $S^1$ factor restores the
$\eta$-like angular mode as a tree-level Goldstone whose mass
arises only from the topological sector, exactly the structure the
dark-energy identification requires.
The machine gauntlet passes Eq.~\eqref{eq:u2pair} \emph{more
cleanly than the $S^3$ model}: particle--hole exchange weights are
zero or repulsive in every channel (the attractive tensor-singlet
of the $S^3$ case cancels away, and the vector-singlet exchange is
repulsive, raising the condensation wall); the exact pairing
decomposition (residual $10^{-14}$) contains only axial diquarks
($c = +0.5$; the $S^3$ model's thinnest channel, the tensor
diquark, and the classic scalar diquark are both identically
absent), with pinned margins $N_{\rm crit} = 0.56$ (gap
convention) and $0.29$ (bubble convention): magnitude-safe at
every $N \geq 1$, with $N=1$ margins $1.8$--$3.5\times$.
Two liabilities are inherited, not created: $\pi_1(U(2)) =
\mathbb{Z}$ restores the axionic-string sector that the $S^1$
model always had (the axion defect-cosmology audit remains an
open dark-energy--paper item), and the Wilson issue is now sharpened from contamination to
absence: at the operating point the pseudoscalar channels are
bubble-supercritical on the Wilson lattice, so the dark-energy
mechanism exists only in a Ginsparg--Wilson discretization---an
existence condition, not an improvement item (quantified in the
dark-energy paper, v5.2).
\textbf{Model ladder verdict: $S^1$ (no matter) $\to$ $S^3$ (no
axion) $\to$ $U(2)$ (both), with the $U(2)$ interaction
Eq.~\eqref{eq:u2pair} the surviving minimal candidate; a $U(3)$
triadic extension (color-like selection structure and baryonic
topology) is recorded as a gated seed in
Section~\ref{sec:u3seed}.}


% ===================================================================
\section{$U(3)$ triadic extension (seed)}
\label{sec:u3seed}
% ===================================================================

\subsection*{Retirement of the old mechanism, with reasons}
An earlier manuscript of this programme (the unified-interactions
draft, 2025) proposed that three flavor-\emph{diagonal} vector
currents $J^{(i)}_\mu = \bar\psi_i\gamma_\mu\psi_i$ generate the
$SU(3)$ algebra through non-vanishing commutators, with
confinement following from the phase constraint
$\sum_i\theta_i = 0$.  That mechanism is retired on documented
grounds: diagonal flavor currents commute, generating only the
Cartan $U(1)^2$, so the claimed algebra was never derived; the
identification of composite vectors with massless gauge fields
was unaddressed (and is now known to face the binding-window and
masslessness obstructions quantified in the vector-channel
paper); and the confinement argument conflated the divergent
energy of an isolated charge with a Wilson-loop area law, without
a mechanism.

\subsection*{The repositioned hypothesis}
What survives is a different and better question, made live by
the present programme's own results: whether the condensate's
\emph{internal order-parameter geometry} extends from $U(2)$ to
\begin{equation}
U(3) \cong \bigl(U(1)\times SU(3)\bigr)/\mathbb{Z}_3,
\qquad \pi_1 = \mathbb{Z},\ \ \pi_2 = 0,\ \ \pi_3 = \mathbb{Z},
\end{equation}
realized microscopically by an $N_f = 3$,
$U(3)\times U(3)$-symmetric chiral interaction (the direct
generalization of Eq.~\eqref{eq:u2pair}).  In this framing the
old triadic-closure intuition acquires a precise and
\emph{structural} meaning: $\sum_i\theta_i = 0$ is the
tracelessness of the $SU(3)$ factor in the $U(3)$ decomposition,
with the singlet (axion) direction carrying the trace---the
constraint is not imposed but automatic.  Baryon number is the
$\pi_3(SU(3)) = \mathbb{Z}$ winding; color-like quantum numbers
label non-singlet order-parameter excitations; and the honest
expectation for ``confinement'' is a \emph{selection structure}
(non-singlet excitations gapped and non-propagating, finite
energy only for singlet windings), not QCD confinement of
constituents.

\subsection*{Gates (pre-registered)}
\begin{enumerate}
\item[0.] \textbf{Realization gate}: the
$U(3)\times U(3)$ interaction at $N_f = 3$ must pass the machine
gauntlet (exact $U(1)_A$ and $SU(3)_A$; particle--hole exchange
zero/repulsive in physical channels; pairing content
absent or protected), as its $U(2)$ predecessor did.  All other
gates are conditional on this one.
\item[1.] \textbf{Topology gate}: verify the homotopy content
above and its compatibility with the $U(2)$ ladder
($\pi_4(SU(3)) = 0$: baryon spin-statistics must then come
entirely from the Wess--Zumino--Witten term).
\item[2.] \textbf{Triadic-closure gate}: confirm the
tracelessness reading against the microscopic model (Cartan
torus vs.\ full non-abelian order); a purely $U(1)^2$ realization
fails this gate.
\item[3.] \textbf{Skyrme/WZW gate (re-scoped v0.20)}: originally
$\kappa_{\rm Skyrme} > 0$; the measured negative $\kappa_U$
re-scopes the stabilization clause to the endogenous-$\omega$
landscape ($R_\star \gg a$), while the WZW/level half is
unaffected.
\item[4.] \textbf{Confinement-structure gate}: a hard criterion
for the selection structure (e.g.\ energetic exclusion of
non-singlet asymptotic excitations), stated so that it can fail;
assertions of area laws are not admissible evidence.
\item[5.] \textbf{$U(2)$-compatibility gate}: the axion
direction, matter topology, pairing safety, and the
chiral-discretization requirement must survive the embedding
unchanged.
\end{enumerate}
\textbf{Gate status (v0.11).}  Gate~0 is \emph{passed}: the
$N_f = 3$, $U(3)\times U(3)$ interaction is machine-verified
exactly $U(1)_A$- and $SU(3)_A$-invariant; its particle--hole
exchange weights are zero in every channel except a repulsive
vector/axial singlet pair (the chiral multiplet receives no shift
at all)---the channel since dynamized as the endogenous
$\omega$ (v0.20 correction note); and its exact pairing decomposition (complete
$66$-channel basis, residual $10^{-14}$) contains only an
axial-diquark flavor-$\mathbf{6}$ at $+1/2$
(magnitude-protected: $N_{\rm crit} = 0.17$--$0.48$ across the
overlap/Wilson bubble conventions, safe at every $N \geq 1$ with
margin $2$--$6\times$) and a \emph{repulsive} vector-diquark
flavor-$\bar{\mathbf{3}}$; scalar-, pseudoscalar-, and
tensor-diquark content is identically absent, so no diquark
condensation route exists.
Gate~1 is bookwork ($\pi_3 = \mathbb{Z}$, $\pi_4(SU(3)) = 0$
placing baryon statistics entirely on the WZW term); Gate~2 is
realized structurally by the $U(3)\times U(3) \to U(3)_V$ order
parameter (tracelessness automatic).  Gate~3's first half is now \emph{passed} by a machine-sealed
embedding argument: the standard $SU(3)$ skyrmion is the Witten
embedding $U = {\rm diag}(U_2, 1)$, whose Maurer--Cartan form has
identically vanishing linear-order projection onto
$\lambda_{4\ldots 8}$ (verified with quadratic-in-$d\phi$
residual scaling) and whose flavor traces coincide exactly with
the $SU(2)$ ones; hence
$\kappa_{\rm Skyrme}^{SU(3),\,{\rm emb}} =
\kappa_{\rm Skyrme}^{SU(2)}$---the embedding identity survives
unchanged, but the inherited value is now the exact
\emph{negative} $-17/(1152\pi^2)$ (v0.20 correction above):
quartic Derrick stabilization is absent, and Gate~3's
stabilization clause is re-scoped to the $\omega$-stabilized
landscape with $R_\star \gg a$ as its quantitative criterion.  The second half---the level---is now anchored by a
nonperturbative machine computation: the spectral flow of the
lattice Dirac Hamiltonian under adiabatic hedgehog creation
(Sylvester-inertia counting of negative levels, Wilson term,
Kahana--Ripka regime $mR \gtrsim 5$) gives net flow $= +1$ for
winding~$1$ (stable on the inter-sector plateau) and $= +2$ for
winding~$2$ at the winding-completion point: the induced baryon
number is winding-proportional with coefficient exactly $1$ per
Dirac species, i.e.\ the Goldstone--Wilczek level is $1$ per
species and the total level is $N$.  (Finite-size threshold
offsets in the crossing locations and the heavy masses used,
$ma = 1.4$--$2.0$, do not affect the integer counting; the full
five-form WZW quantization and the spin-statistics linkage then
follow by the standard consistency argument, cited rather than
machine-verified.)  The stake is thereby placed on computed
ground: \emph{baryons are fermionic iff $N$ is odd}---the species
number is no longer a free parameter but is constrained by
baryonic statistics.
Gate~4 is now resolved by a three-tier criterion with honest
verdicts (v0.15).
\emph{Gate 4A-weak: PASS} --- open (non-closed) color-phase
configurations carry angular gradients $\partial U \sim 1/r$ at
infinity and hence linearly divergent energy, so only closed,
constant-at-infinity configurations are finite-energy matter
candidates; combined with $\pi_3 = \mathbb{Z}$,
$\kappa_{\rm Skyrme} > 0$, and level $= N$, closed baryonic
winding is finite-energy and topologically protected.
\emph{Gate 4A-strong: NOT PASSED} --- finite energy does not
exclude non-singlet global orientations.
\emph{Gate 4B: FAIL for the global theory} --- collective
quantization of the skyrmion's $SU(3)$ orientation produces a
rotor spectrum of finite-energy multiplets
($E \sim M + C_2(R)/2\mathcal{I}$), so the physical Hilbert
space is not singlet-only.  The correct reading is
classificatory: the $U(3)$ seed as it stands is a
\emph{flavor-like} structure delivering baryonic topology and
multiplet spectroscopy; ``color confinement'' is removed from
its claims, and its possible restoration is recorded as a fork
with stated costs (fundamental gauge fields abandon the
substrate-only premise; emergent gauge redundancy faces the
vector-channel binding obstruction).
\emph{Gate 4C} (flux/area law) is deferred as a future stronger
test, meaningful only past the fork.
One productive corollary of the 4B failure: the WZW term
constrains the allowed rotor representations
(Guadagnini-type quantization, reps containing hypercharge
$N/3$), so the multiplet content is an $N$-dependent
\emph{prediction}: $N = 3$ selects octet/decuplet-like
spectroscopy while also satisfying the odd-$N$ fermionic-baryon
constraint---the second independent place the species number is
physically pinned.  The quantization-condition check is now complete
(machine-verified enumeration): with level $k = N$ from the
spectral-flow computation, the Guadagnini--Witten constraint
$Y_R = N/3$ selects minimal representations $p + 2q = N$, giving
$\mathbf{3}_{1/2}$ at $N{=}1$;
$\bar{\mathbf{3}}_0 \oplus \mathbf{6}_1$ at $N{=}2$; and
$\mathbf{8}_{1/2} \oplus \mathbf{10}_{3/2}$ at $N{=}3$---the
QCD-like octet/decuplet appearing first and only at $N = 3$.
An internal consistency check passes automatically:
$p \equiv N \pmod 2$ forces half-integer rotor spin exactly when
$N$ is odd, so the spectroscopic and statistics constraints on
$N$ are redundantly coherent.  \textbf{The species number is
thereby pinned from three directions: WZW level, fermionic
baryon statistics, and octet/decuplet spectroscopy---all
selecting $N = 3$.}
Caveats, attached permanently: the rigid-rotor quantization
assumes finite moments of inertia from the induced Skyrme+WZW
Lagrangian (not yet computed); the selection is conditional---%
\emph{if} the substrate's baryonic windings are to reproduce
observed baryon spectroscopy, then $N = 3$---and the
identification with physical baryons requires the
Standard-Model-embedding step that leg~(1) of the
Lorentz-invariance trichotomy independently demands.
The $U(3)$ seed's mature claim is accordingly: \emph{baryonic
topology with $N = 3$ selection}, not confinement.
\emph{Spectral-test anchor (v0.14).}  The overlap $\sigma$-pole
scan across $m_{\rm dyn} a = 0.1$--$0.4$ measures
$m_\sigma/(2m_{\rm dyn}) = 1 - 0.98\,(m_{\rm dyn}a)
+ 0.72\,(m_{\rm dyn}a)^2$ (five points, residuals
$\leq 0.007$, volume systematic $0.7\%$): the chiral-limit
intercept is unity within errors, recovering the continuum
relation $m_\sigma = 2m_{\rm dyn}$ exactly, with the lattice
correction now a measured law.  Together with the exact Goldstone
$m_\pi = 0$, the mass side of the spectral test of
Section~\ref{sec:test} is fully anchored; the coupling gate
remains its outstanding requirement.
Verdict discipline: this section is a seed; no companion paper is
opened until Gates 3--4 also have first results.


% ===================================================================
\section{The $S^2$ layer: orientation sector and emergent
$U(1)$ connection (gated seed)}
\label{sec:s2seed}
% ===================================================================

The Hopf fibration $S^1 \to S^3 \to S^2$ identifies a middle layer
of the ladder that is \emph{not a new rung}: by the Vafa--Witten
theorem, vector symmetries of the substrate do not break
spontaneously, so $S^2 = SU(2)/U(1)$ is unavailable as a vacuum
manifold of a new condensate; instead, $S^2 \cong \mathbb{CP}^1$
is the \emph{directional quotient of the already-adopted $U(2)$
model's chiral field}---writing the $SU(2)$ sector as a normalized
spinor $z$ with $\vec n = z^\dagger\vec\sigma z$, the layer comes
with a composite Berry connection
$a_\mu = -i z^\dagger\partial_\mu z$,
$f_{\mu\nu} \sim \vec n\cdot(\partial_\mu\vec n\times
\partial_\nu\vec n)$, at no cost in new degrees of freedom.
Its honest physics menu: bosonic orientation modes; monopole
($\pi_2 = \mathbb{Z}$) and Hopfion ($\pi_3(S^2) = \mathbb{Z}$)
topology; and a photon-\emph{candidate} connection that is
emphatically not yet a photon.

Gates, pre-registered:
\begin{enumerate}
\item[S2-1.] \textbf{Induced kinetic term: FAILED as originally
stated (v0.20).}  The inheritance identity itself is intact and
machine-exact ($[L_\mu,L_\nu] = -2i(\vec n\cdot\vec\tau)
f_{\mu\nu}$, so the layer-restricted Skyrme term \emph{is} the
Berry--Maxwell term), but the inherited sign has flipped: with
the corrected $\kappa_U = -17/(1152\pi^2) < 0$ the induced
Berry--Maxwell coefficient is \emph{negative}---a wrong-sign
kinetic term for the photon candidate at one loop.  The
Vakulenko--Kapitanskii Hopfion-stabilization bonus is withdrawn
with it.  Recorded outs, uncomputed: contributions to the Berry
kinetic term beyond the layer-restricted quartic; and the
demonstrated fact that not every composite vector is poisoned
(the isoscalar $\bar\psi\gamma^\mu\psi$ channel's induced
kinetic term is healthy, $Z_V > 0$---the ghost is specific to
the $\mathbb{CP}^1$ Berry connection).  Gate S2-2's
Coulomb-phase framing is moot in its current form until S2-1 is
resolved.
A framing correction is recorded alongside: masslessness in this
programme has three distinct protections---Goldstone (the
$\vec\pi,\eta$ modes), gauge redundancy (the $a_\mu$
candidate, whose only breach is the monopole gate S2-2), and
none for solitons, whose masses \emph{are} their topological
energies; ``$S^2$ massless vs $S^3$ massive'' is a category
comparison between a connection and a soliton, not between
manifolds.
\item[S2-2.] \textbf{The monopole gate (central): restated
(v0.22), and its Coulomb-side resolution retracted.}
The gate's earlier statement---that in $3{+}1$D a compact
emergent $U(1)$ with dynamical monopoles is ``generically gapped
by the Polyakov-class mechanism''---was wrong, and is corrected
here following the companion dark-energy paper's round-one
review.  Polyakov's monopole-plasma gapping is a $2{+}1$D
(three-dimensional Euclidean) result; in $3{+}1$D compact $U(1)$
possesses a genuine \emph{Coulomb phase} at weak coupling, with
an \emph{exactly} massless photon protected by the phase itself
(Guth), and monopole condensation---hence gapping and
confinement---occurs only on the strong-coupling side of a phase
transition.  The gate is therefore not ``exhibit a suppression
mechanism or lose the photon'' but a \emph{phase question}: does
the emergent compact $U(1)$ of this layer sit on the Coulomb or
the confining side?
The answer that v5.6--v5.8 of the companion paper recorded on the
Coulomb side is \textbf{retracted}: it rested on the superseded
(fit-leakage) $\kappa$ numbers of the July-8 run.  With the
corrected coefficient the relevant arithmetic runs the other way.
The same monopole-loop action that the companion paper's
dark-energy chain required, expressed through the measured
coefficient, is
\begin{equation}
S_{\rm mono} \;=\; 640\,\kappa_U \;=\; -\frac{85}{9\pi^2}
\;=\; -0.95 ,
\end{equation}
against a required window $[140, 550]$: the pre-registered
criterion fired, and the companion paper records the termination
of that chain.  The reading relevant \emph{here} is the physical
one and survives the sign: the substrate's quartic resistance to
orientation twisting is of ordinary one-loop (Gasser--Leutwyler)
strength, $|S_{\rm mono}| \sim 1$, so monopoles in this layer are
\emph{cheap} and topological events dense rather than
exponentially rare.  This sector supplies no diluteness of its
own; a Coulomb-phase claim would have to come from elsewhere.
All of which is, in the present state, \emph{moot}: the
Coulomb-versus-confined question presupposes a Maxwell term whose
sign is right, and S2-1 reports that at one loop it is not.  The
gate is suspended behind S2-1 rather than answered, and remains
stated so that it can fail.
\item[S2-3.] \textbf{The $\eta$--Hopf coupling (dark-energy
payoff)}: whether the fermion determinant couples the angular mode
to the Hopf density $f\wedge f$ of the Berry connection is a
well-posed derivative-expansion question; if yes, Hopfion
topological fluctuations of the orientation sector constitute a
\emph{third} candidate source for the dark-energy paper's
re-scoped topological susceptibility, distinct from both the
(disfavored) channel-composite gauge sector and the gravitational
$R\tilde R$ route.
\item[S2-4.] \textbf{Compatibility}: the layer must not disturb
the $U(2)$ model's verified protections (it introduces no new
operators, so this is expected but must be checked where the
$\mathbb{CP}^1$ decomposition interacts with the Wilson/overlap
discretization).
\end{enumerate}
No claims are attached to this section beyond the existence of the
decomposition; in particular ``$S^2$ = photon'' is explicitly not
asserted, and the layer's value is conditional on Gates S2-1
through S2-3.

% ===================================================================
\section{Gates and next actions}
\label{sec:gates}
% ===================================================================

\begin{enumerate}
\item \textbf{Fierz arithmetic: COMPLETED (v0.19 status), with the
opposite of the anticipated outcome.}  The machine-verified
decomposition of the minimal chiral pair yields identically zero
vector and axial couplings by Fierz exchange (and a repulsive
tensor), so the cross-channel ratios $G_V/G_S$, $G_A/G_S$ are
\emph{not} Fierz-inherited; in the adopted invariant-universal
model they are fixed by the single-coupling economy postulate
instead ($G_V = G_A = G_\chi$), which is a model assumption, not
an output.  The Aoki-risk sub-question (b) is resolved by the
frame requirement: the operating-point analysis must be performed
in the overlap discretization, where the $S/P$ sector is exactly
degenerate and healthy.
\item \textbf{Chirally clean scalar extraction: COMPLETED.}
Overlap true-pole scan; $m_\sigma$ law measured with chiral-limit
intercept $1.000(10)$ (Sections~\ref{sec:partial}
and~\ref{sec:u3seed}).
\item \textbf{Ratio-fit design (re-scoped).}  Before the
$N \le 5$ fit of Section~\ref{sec:test} can run, the state list
must be settled under the invariant-universal coupling and the
vector-state caveat: at $G = G_\chi$ the vector channel has no
below-threshold state (companion vector paper), so the candidate
tower presently consists of $m_\pi = 0$, $m_\sigma$, and
whatever the $V/A$ sector contributes on the resonance side.
Whether a two-to-three--state tower constitutes a meaningful test
(kill criterion (ii) of Section~\ref{sec:test} is close to
triggering) is the design question to resolve first; if it
triggers, the negative result is recorded per the version
discipline below.
\item \textbf{Composite--composite cone precision campaign}
(unchanged; programme-level priority, independent of the tower
test): resolve the $P$--$V$ split scaling law below the current
systematic floor (Section~\ref{sec:epgeom}) at the interacting
level.  This gates leg~(2) of the survival trichotomy and
therefore the viability of the entire programme against
Lorentz-invariance bounds.
\end{enumerate}

\workingnote{The dark-energy scale $\Lambda_{\mathrm{DE}}$ is
deliberately excluded from the tower: it is exponentially
suppressed instanton physics ($\Lambda^4 e^{-S}$), a different
mechanism from any $1/(Nb)$ spectrum, and no small-$N$ tower spans
thirty decades. Do not let future drafts re-import it.}

\workingnote{Version discipline: this note becomes ``Paper 5
v1.0'' only if and when the spectral test runs and passes.
If the test fails, the negative result is still worth a short
note --- constraint on the internal structure of the condensate
--- consistent with the programme's honesty-first policy.}

\end{document}
